\section{A second construction using remote odd powers}\label{sec:remote}
A second profile construction gains fourteen cycles for every four
degrees. Each stationary inflection splits on both reflected sides.
Moment conditions at separated radii prescribe two cubic fold windows
at each inflection. A gap between the highest odd power and the even
degree allows further sign changes on the last rising branch while
the earlier samples stay fixed. Propositions~\ref{prop:physical-quad},
\ref{prop:cap-transfer} and~\ref{prop:sard} realize this single profile
and close its orbits. The resulting bounds have the same asymptotic
coefficient as Theorem~\ref{thm:main-asymptotic} and smaller finite-degree
constants.

\begin{proposition}[Cycle counts from remote powers]\label{prop:remote-odd}
For integers $h\ge1$ and $e\in\{0,1,2,3\}$ one has
\begin{equation}\label{eq:remote-formula}
 \HH(10+4h+e)\ge14h+c_e,\qquad(c_0,c_1,c_2,c_3)=(9,10,13,14).
\end{equation}
The bounds hold on nonempty open sets of exact-degree coefficients.
In particular $\HH(n)\ge(7n-63)/2$ for $n\ge14$.
\end{proposition}

\subsection{Two cubic windows at separated heights}
For real $a<0,b$, let $f(v)=v^3+av+b$, $k=\sqrt{-a/3}$,
$\mu=b-2k^3$, $M=b+2k^3$, and
\[
 P_a(v)=\int_0^v(3r^2+a)^2\,\dd r=\tfrac95v^5+2av^3+a^2v.
\]
The largest increasing inverse skips $[-2k,k]$ at $\mu$.
For $T\ne\mu,M$ let $v$ be the largest root of $f(v)=T$ and set
\begin{align}
 J^I_{a,b}(T)&=P_a(v)-\ind_{T>\mu}\{P_a(k)-P_a(-2k)\},\nonumber\\
 J^U_{a,b}(T)&=
 \begin{cases} P_a(-2k),&\mu<T<M,\\ J^I_{a,b}(T),&\text{otherwise}.
 \end{cases}                                                     \label{eq:remote-cubic}
\end{align}
The half-returns follow the original flow forward on the right and
backward on the left. In this orientation the physical ordinate, hence
the cubic height, decreases. An entrance from the right follows the
largest root in $J^I_{a,b}$. For a launch with $\mu<T<M$, the singular
launch path follows the inner rising branch. The reference follows the
outer rising branch to $v=k$ and jumps to $v=-2k$ at height $\mu$.
Just below $\mu$ both singular paths follow the inner branch with $v<-2k$.
Thus $-2k$ is the upper integration endpoint for their common segment
in $J^U_{a,b}(T)$.

\begin{lemma}[Action differences at two cubic windows]\label{lem:skew-cubic}
For every sufficiently large $C>0$, there exist $\alpha<-1$ and $\Lambda$
for which the two differences
\[
 D^X(T)-\Lambda
 =J^X_{\alpha+1,-C}(T)-J^X_{\alpha-1,C}(T)-\Lambda,\qquad X=U,I,
\]
have, at finite regular samples in increasing order, respectively seven
and three strict changes. Both lists start positive and end negative.
The samples for the two lists can be chosen independently.
\end{lemma}
\begin{proof}
Write $a_\pm=\alpha\pm1$, $k_\pm=\sqrt{-a_\pm/3}$. The windows will be
disjoint. At $M_+=-C+2k_+^3$ the upper difference jumps up by
$P_{a_+}(2k_+)-P_{a_+}(k_+)>0$. At $M_-=C+2k_-^3$ it jumps down by
the corresponding positive amount for the minus profile. We equalize the
midpoints $L_C,R_C$ of the two jumps.

Take $\alpha=-t C^{2/9}$ with $1/4\le t\le3$. For large $C$ both
$a_\pm<0$, and the two window widths are $O(C^{1/3})$. At $Y\gg|a|^{3/2}$
the positive root $v^3+av=Y$ satisfies
\begin{equation}
 P_a(v)=\tfrac95Y^{5/3}-aY+O(a^2Y^{1/3}).                     \label{eq:inverse}
\end{equation}
Indeed $v/Y^{1/3}=1-a/(3Y^{2/3})+O(a^2/Y^{4/3})$ by the implicit
function theorem at 1. Oddness gives the negative-root expansion at
large negative $Y$. At the first jump the regular minus root has
$|Y|=2C-2k_+^3$. At the second the regular plus root has
$Y=2C+2k_-^3$. All the endpoint and skip terms involving only $k_\pm$
are $O(k_+^5+k_-^5)=o(C)$. Expansion of \eqref{eq:inverse} at $2C$
therefore gives, uniformly in $t$,
\[
 \frac{R_C-L_C}{C}=
 12\,2^{2/3}(t/3)^{3/2}-4+o(1).
\]
The first term comes from the sum $2(k_-^3+k_+^3)$ of the height
increments and the second from $a_+-a_-=2$. This expression has opposite
signs at $t=1/4$ and $t=3$. Continuity in $t$ gives a common midpoint
$\Lambda\sim(9/5)(2C)^{5/3}$.

At $T=0$ both differences are asymptotic to $(18/5)C^{5/3}$. Their
ratios to $\Lambda$ tend to $2^{-2/3}<1$.
At $T=C^2$ both are asymptotic to $6CT^{2/3}$, by
\eqref{eq:inverse} at $T+C,T-C$. Here the $O(T)$ linear term, the
$O(|\alpha|^2T^{1/3})$ error and the skips are all lower order.
Thus both exceed $\Lambda$ there. With $C,\alpha$ now fixed,
$D^U,D^I=-2T+O_{C,\alpha}(|T|^{2/3})$ at either infinity.
Take successive samples far to the left, immediately before and after
$M_+$, at 0, immediately before and after $M_-$, at $C^2$, and far
to the right. These give $+,-,+,-,+,-,+,-$ for the upper difference.
The two tails, 0 and $C^2$ give $+,-,+,-$ for the incoming difference.
All samples stay off both pairs of critical values.
\end{proof}

\subsection{Simultaneous preparation at separated radii}
Write $u=z^2/2$ and fix $h\ge1$. Here $B_u$ denotes the derivative of $B(\sqrt{2u})$. Choose separated $0<q_1<\cdots<q_h$
and set
\[
 B(z)=\tfrac32 z^3\prod_{i=1}^h(1-z^2/q_i^2),\qquad
 E_0(z)=\int_0^z vB(v)^2\,\dd v,\qquad J=3+2h .
\]
We call each $q_i$ a \emph{shoulder} of the unperturbed profile.
It is a stationary inflection. The polynomial $B$ has a simple zero there, so $E_0'=zB^2$
has a double zero and $E_0'''(q_i)=2q_i B'(q_i)^2>0$. The prepared
perturbation opens a maximum and a minimum near $q_i$ on each side.
The initial even degree is $m_0=2J+2=8+4h$. Write $s_i=(-1)^i$.

\begin{lemma}[Simultaneous sign lists at separated radii]\label{lem:remote-base}
There are positive profiles $G_s=E_\eta+sC$ of the above type, of even degree
$m_0$ and highest odd degree $j_0=J+2$, with germs
\[
 E_\eta(z)=\eta z^2+O_\eta(z^8),\qquad
 C(z)=\tfrac32\eta z^3+O_\eta(z^5)
\]
for a fixed $\eta>0$. At each shoulder $q_i$ of $E_0$, in increasing order of
positive regular height, the signs of $J_+^U-J_-^U$ and $J_-^I-J_+^I$ are
\[
 s_i(+,-,+,-,+,-,+,-),\qquad -s_i(+,-,+,-),
\]
respectively. At increasing negative entrance magnitudes $w$ there are
$h+1$ alternating signs of $N(w)$ beginning negative. The leading sign
of $C$ is $\beta_0=-(-1)^h$.
\end{lemma}
\begin{proof}
Choose the radii and the odd height-splitting mode $H$ together. Write $H=z^5C_0(z^2)$ with $C_0(0)=-1$ and $\deg C_0\le h$. Require
\[
 \int_0^{q_i}B^2H'\,\dd z=0,\qquad 1\le i\le h.
\]
The prefix
integral equations are square in the remaining coefficients. Subtracting
successive equations gives one zero integral on each $(q_{i-1},q_i)$,
$q_0=0$. In the homogeneous system $H_z/z^4$, a polynomial in
$z^2$ of degree at most $h$, would have a positive root in each interval
and a root at zero. Hence the system is invertible. For its solution
the same $h$ interior roots force the top degree.

The case of zero shoulders has $B=(3/2)z^3$, $H=-z^5$.
Assume the first $i$ radii and their moment solution have been chosen. Suppose the old top terms are $b z^j$, $a z^{j+2}$, with $ab<0$.
Append $q=L$ by multiplying $B$ by $1-z^2/L^2$. In the new system
write the odd shape as an old-degree trial polynomial plus
$\lambda L^{-2}z^{j+4}$. The old rows converge to the preceding
invertible system, with zero limiting coefficient in the new column.
In the new row put $z=Lv$ and divide by $b^2 L^{3j+2}$. In the limit
the old rows have fixed the trial top coefficient at $a$, and the new
row is
\[
 \int_0^1 v^{2j}(1-v^2)^2\big((j+2)a v^{j+1}
                            +(j+4)\lambda v^{j+3}\big)\,\dd v=0.
\]
The new diagonal coefficient is positive. Thus all old coefficients
converge and $\lambda\to-aD_j$, where
\[
 D_j=\frac{(j+2)(3j+8)}{(j+4)(3j+2)}>1 .
\]
The numerator exceeds the denominator by $8$.
Consequently $H(L)/L^{j+2}\to a(1-D_j)$, of the opposite sign to $a$.
The sign of $B_z(L)$ is opposite to $b$, and the new top signs are
again opposite. Strict inequalities at earlier shoulders persist. All
limits here follow coefficientwise after the indicated substitution.

For $w>0$ write
\[
 K(w)=\int_0^\infty \frac{zB(z)}{(w+E_0(z))^2}\,\dd z .
\]
Fix the old flight samples before enlarging $L$. The new primitives and
odd shapes converge on compact sets. For large $z$ the contribution of
the old top square to the new $E_0$ equals
\[
 b^2 z^{2j+2}\left(\frac1{2j+2}-\frac{2(z/L)^2}{2j+4}
                                +\frac{(z/L)^4}{2j+6}\right).
\]
The bracket is bounded below by a positive multiple of $1+(z/L)^4$,
as follows by integrating a square or by its negative discriminant.
Every lower contribution is $O(z^{2j}(1+(z/L)^4))$ for $z\ge1$,
uniformly in $L\ge1$. The new $|B|$ is bounded by
$Cz^j(1+(z/L)^2)$. These give an integrable tail bound for every old
sample. At large $w$ the new $K$ has the sign of the new top odd
coefficient. Scaling $z=w^{1/\deg E_0}v$ in the top-degree range gives
a nonzero signed limit after division by $w^{(\deg B+2)/\deg E_0-2}$.
Bounded abscissas contribute $O(w^{-2})$. Thus the new large-$w$ sign is that of $-b$, whereas the preceding
large-$w$ sign was that of $b$. At the zero-shoulder start,
$B=(3/2)z^3$ makes the integrand of $K$ positive for every $w>0$.
Keeping each old sample before taking the next large $w$ proves the induction.

At each radius put $u_i=q_i^2/2$ and set
$b_i=|B_u(u_i)|$, $l_i=\sqrt{3/b_i}$, $d_i=b_i l_i$ and $r_i=b_i^2 l_i$.
Choose $g>0$ so that $gH(q_i)/(s_i d_i)=-C_i$ is negative of sufficiently
large magnitude at every $i$. The cubic comparison in
Lemma~\ref{lem:skew-cubic} gives $\alpha_i,\Lambda_i$ for the two graphs
$v^3+(\alpha_i\pm1)v\mp C_i$ and the two lists in the statement.
To specify their offsets on the polynomial, put
$a_{s,i}=\alpha_i+s s_i$. For $a<0$ let
\[
 p_a(v)=9v^5/5+2av^3+a^2v,\quad k=\sqrt{-a/3},\quad
 d(a)=p_a(k)-p_a(-2k).
\]

For the two even controls consider the map
\[
 {\cal L}(V)=\left((V(q_i))_{i=1}^h,
          \left(\int_0^{q_i}V B'\,\dd z\right)_{i=1}^h\right),
 \qquad V\in z^6\R_{\le 2h}[z^2].
\]
If $\mathcal L(V)=0$, divisibility at the $q_i$ and evenness give $V=BW$
with $W=z^3P(z^2)$, $\deg P\le h$. Consecutive moment differences and
$B(q_i)=B(0)=0$ give $\int_{q_{i-1}}^{q_i}B^2 W'\,\dd z=0$, $q_0=0$.
If $[z^6]V=0$ then $W=z^5P_0(z^2)$, $\deg P_0\le h-1$.
Consequently $W'/z^2$ has degree at most $h$ in $z^2$ and vanishes at zero.
Unless identically zero it also has a positive root in each of the $h$
disjoint intervals, a contradiction. Thus $W$ is zero. The restriction of $\mathcal L$ to
$\mathcal V_0=z^8\R_{\le 2h-1}[z^2]$ (zero $z^6$ coefficient) is
therefore an isomorphism onto $\R^{2h}$. The full space has dimension
$2h+1$ and its one-dimensional kernel is spanned by $B^2$, since
$\int_0^{q_i}B^2B'\,\dd z=0$. Choose $D=2+V$ with $V\in\mathcal V_0$
first, and then $Z\in\mathcal V_0$, by specifying all $2h$ coordinates
of each image:
\begin{align}
 D(q_i)&=\alpha_i b_i,& \int_0^{q_i} D B'\,\dd z&=0,\nonumber\\
 Z(q_i)&=0,&4\int_0^{q_i}(Z B'+g D H')\,\dd z
 &=\sum_{l<i}r_l\{d(a_{+,l})-d(a_{-,l})\}-s_i r_i\Lambda_i .
                                                              \label{eq:offset}
\end{align}
The constant two in $D$ supplies the quadratic germ. Its contribution
to the last row, including $2gH(q_i)$ inside the integral value, is
included when choosing $Z$.

Set $(E_\eta)_u=B^2+\eta D+\eta^{3/2}Z$, $E_\eta(0)=0$, and
$C=\eta B+g\eta^{3/2}H$. For small $z$ both $u$-slopes dominate a
positive multiple of $z^6+\eta$, since $B_u=O(z)$. At the other end
$B^2$ dominates the perturbations in leading coefficient. On intermediate
compact sets outside the roots it gives a positive slope as well. Near
$u_i$, in $u=u_i+l_i\sqrt\eta\,v$ and
$t=E_\eta(q_i)+d_i\eta^{3/2}T$, the two profiles in $T$ converge in $C^3$
on bounded $v$ intervals to
$v^3+a_{s,i}v-s s_i C_i$. Outside a large fixed multiple of $\sqrt\eta$
the quadratic growth of $B^2$ in $u-u_i$ dominates the perturbation of
the slope. This gives just two simple folds there on each side.
Their heights tend to distinct positive values for distinct $i$.
In particular the profiles are positive away from zero.

Before skips, the difference of actions up to $u_i$ is
$4\int_0^{u_i} (E_\eta)_u C_u\,\dd u$. Its terms of orders $\eta,\eta^{3/2},\eta^2$
vanish by $\int B^2 B'\,\dd z=B^3/3$ at the endpoints, the $H$ moments and
\eqref{eq:offset}. The remaining leading term is the left side of
\eqref{eq:offset} times $\eta^{5/2}$. The local primitive is
$r_i\eta^{5/2}p_{a_{s,i}}(v)+o(\eta^{5/2})$. Each earlier completed
jump deletes $r_l\eta^{5/2}d(a_{s,l})+o(\eta^{5/2})$. An axis launch
trapped inside a window uses the landing $v=-2k$, while an incoming
launch uses the largest entry root and the deletion at its minimum.
Thus \eqref{eq:offset} gives exactly the common offset of the cubic
comparison, with the physical side labels interchanged for $s_i=-1$.
All selected inequalities are strict for one sufficiently small $\eta$.
Finally $N(w)/(-2\eta)\to K(w)$ at all the finite negative samples by
dominated convergence with the weight $1+z^{m_0}$. This proves the lemma.
\end{proof}

\subsection{Adding the odd powers and closing the incoming list}
\begin{proof}[Proof of Proposition~\ref{prop:remote-odd}]
At the last shoulder of Lemma~\ref{lem:remote-base} the upper,
positive-incoming and negative-incoming lists end with signs
$(\beta_0,-\beta_0,\beta_0)$. Lemma~\ref{lem:tail-action} at
$j_0=5+2h$ gives the opposite signs farther out. Now introduce the
$h$ odd powers of degrees $j_0+2,\ldots,m_0-3$ in order. Choose each
coefficient with sign opposite to the preceding leading odd term and
small enough to preserve every sample already fixed. Only then choose
the new positive and negative ordinates far enough out to record the
next signs. Each list gains $h+1$ changes, and the final leading odd
coefficient has sign $(-1)^h\beta_0=-1$.

For $e\in\{0,1,2,3\}$ take the first $e$ terms of the ordered list
\[
  \bigl(-\delta_1 z^{m_0-1},\quad \delta_2 z^{m_0+2},\quad
        -\delta_3 z^{m_0+1}\bigr).
\]
At each step choose $\delta_i>0$ small enough to preserve the previous
samples, then select the new ordinates. Each addition switches
the relation between the current odd and even degrees from $j<m-2$
to $j=m-1$ or conversely. This adds a change to the upper and
positive-incoming lists and leaves the last negative-entrance sign
positive. With all ordinates fixed, choose the common reference and
apply Proposition~\ref{prop:physical-quad}.

Proposition~\ref{prop:physical-quad}(i) gives $\D(h_*)>0$ for a
sufficiently small $h_*>0$, preceding the negative sign at the
first shoulder. The shoulders contribute $7h$ internal changes,
and the last sign of each agrees with the first of the next. The far levels
contribute $h+1+e$ changes. Detuning the reference to the sign opposite
the last sample adds the final change. Proposition~\ref{prop:sard}(i)
then gives one linear coefficient with a simple zero in every selected
bracket, hence $q=1+7h+(h+1+e)+1=8h+3+e$ distinct compact hyperbolic cycles.

In decreasing ordinate the positive incoming list has $3h+(h+1)+e$
changes and ends positive. Its join to the first negative-ordinate
sign contributes one. The negative list has $h+(h+1)$ changes.
Thus the unshifted total is $6h+3+e$ and the final sign is positive.
For even $e$, take one sufficiently small positive $\Delta$ preserving
the whole list, then append a sufficiently negative ordinate. Indeed for
the fixed physical polynomial of even degree $m$, $I_s(-w)=w+O(w^{2/m-1})$
as $w\to\infty$ by the same positive integral equation. The translated
difference tends to $-\Delta$. This supplies one further change.
For odd $e$ take $\Delta=0$.

Use the matched-maximum theorem of Proposition~\ref{prop:cap-transfer}. For an
input of even degree $m$ and highest odd degree $j=m-g$, $g=1$ or $3$,
with top odd coefficient $b$, it gives
\[
 q+r+1+\mathbf1_{\{\sigma\ne\sgn((2-g)b)\}}
 \quad\text{at exact degree }\max(m+2,j+4).
\]
Here $r$ counts incoming changes at a single translation and $\sigma$
is the final sign. In this construction $b<0$. The final sign is
negative for $g=3$ and positive for $g=1$, so the
indicator is always one. The exact degrees and counts are
\[
\begin{array}{c|cc|cc|cc}
 e & m & j & q & r & \deg F_{\rm final}&\text{cycles}\\
 \hline
 0&m_0&m_0-3&8h+3&6h+4&m_0+2&14h+9\\
 1&m_0&m_0-1&8h+4&6h+4&m_0+3&14h+10\\
 2&m_0+2&m_0-1&8h+5&6h+6&m_0+4&14h+13\\
 3&m_0+2&m_0+1&8h+6&6h+6&m_0+5&14h+14 .
\end{array}
\]
Hyperbolic persistence gives the asserted open sets at each degree.
\end{proof}
